\documentclass[10pt,letterpaper]{article}

% Typography
\usepackage{lmodern}
\usepackage{tgpagella}
\usepackage[T1]{fontenc}
\usepackage[utf8]{inputenc}
\usepackage{microtype}
\usepackage{amsmath,amssymb,amsthm,mathtools,bm}
\usepackage{siunitx}

% Layout
\usepackage[margin=0.75in,top=0.62in,bottom=0.62in]{geometry}
\usepackage{parskip}
\usepackage{setspace}
\usepackage{graphicx}
\usepackage[dvipsnames,table]{xcolor}
\usepackage{tikz}
\usetikzlibrary{shapes,decorations,shadows}

% Tables and lists
\usepackage{array}
\usepackage{tabularx}
\usepackage{booktabs}
\usepackage{longtable}
\usepackage{enumitem}
\usepackage{multicol}
\usepackage{pbox}

% Navigation and styling
\usepackage{fancyhdr}
\usepackage{lastpage}
\usepackage{titlesec}
\usepackage{fontawesome5}
\usepackage{hyperref}
\usepackage{tcolorbox}
\tcbuselibrary{skins,breakable}

\definecolor{primaryblue}{RGB}{0,47,108}
\definecolor{accentblue}{RGB}{44,130,201}
\definecolor{darkgray}{RGB}{64,64,64}
\definecolor{lightgray}{RGB}{245,245,245}
\definecolor{urgold}{RGB}{255,204,0}

\hypersetup{
    colorlinks=true,
    linkcolor=accentblue,
    urlcolor=accentblue,
    citecolor=primaryblue,
    pdftitle={LING 414: Final Project Rubric},
    pdfauthor={Aaron Steven White}
}

\renewcommand{\familydefault}{\sfdefault}
\setstretch{1.02}
\setlist[itemize]{leftmargin=*,itemsep=1pt,topsep=2pt,parsep=0pt}
\setlist[enumerate]{leftmargin=*,itemsep=1pt,topsep=2pt,parsep=0pt}

\titleformat{\section}
  {\large\bfseries\color{primaryblue}}
  {\thesection}{0.7em}{}
  [\vspace{-0.25em}\color{accentblue!55}\titlerule]
\titlespacing*{\section}{0pt}{0.8em}{0.35em}

\titleformat{\subsection}
  {\normalsize\bfseries\color{darkgray}}
  {\thesubsection}{0.7em}{}
\titlespacing*{\subsection}{0pt}{0.65em}{0.25em}

\titleformat{\subsubsection}
  {\small\bfseries\color{darkgray}}
  {\thesubsubsection}{0.7em}{}
\titlespacing*{\subsubsection}{0pt}{0.5em}{0.2em}

\pagestyle{fancy}
\fancyhf{}
\fancyhead[L]{\scriptsize\color{darkgray}LING 414 | Fall 2026}
\fancyhead[R]{\scriptsize\color{darkgray}Final Project Rubric}
\fancyfoot[C]{\scriptsize\color{darkgray}Page \thepage\ of \pageref{LastPage}}
\renewcommand{\headrulewidth}{0.4pt}
\renewcommand{\footrulewidth}{0pt}

\newtcolorbox{bluebox}[1][]{
  colback=accentblue!8,
  colframe=accentblue!65,
  boxrule=0.8pt,
  arc=2mm,
  left=7pt,right=7pt,top=5pt,bottom=5pt,
  #1
}

\newtcolorbox{graybox}[1][]{
  colback=lightgray,
  colframe=darkgray!65,
  boxrule=0.7pt,
  arc=2mm,
  left=7pt,right=7pt,top=5pt,bottom=5pt,
  #1
}

\newtcolorbox{goldbox}[1][]{
  colback=urgold!12,
  colframe=urgold!90!black,
  boxrule=0.8pt,
  arc=2mm,
  left=7pt,right=7pt,top=5pt,bottom=5pt,
  #1
}

\newtcolorbox{redbox}[1][]{
  colback=red!7,
  colframe=red!55,
  boxrule=0.8pt,
  arc=2mm,
  left=7pt,right=7pt,top=5pt,bottom=5pt,
  #1
}

\begin{document}

\begin{center}
\colorbox{primaryblue}{\parbox{0.94\textwidth}{%
\vspace{0.34cm}
\begin{center}
  \color{white}
  {\LARGE\bfseries Final Project Rubric}\\[0.18em]
  {\normalsize LING 414: Statistical Methods in Linguistics}\\[0.15em]
  {\normalsize Fall 2026}\\[0.25em]
  \rule{0.42\textwidth}{0.4pt}\\[0.18em]
  {\small University of Rochester}\\
  {\small Department of Linguistics}
\end{center}
\vspace{0.34cm}
}}
\end{center}

\section{Overview}

\begin{bluebox}
\faInfoCircle\ \textbf{Final Project Weight:} 40\% of total course grade (LING 414)
\begin{itemize}
  \item \textbf{Pre-proposal Meeting:} 5\% of project grade (2\% of course grade)
  \item \textbf{Proposal:} 5\% of project grade (2\% of course grade)
  \item \textbf{Presentation:} 20\% of project grade (8\% of course grade)
  \item \textbf{Code Submission:} 20\% of project grade (8\% of course grade)
  \item \textbf{Paper Submission:} 50\% of project grade (20\% of course grade)
\end{itemize}
\end{bluebox}

\section{Project Options}

Students must choose one of three project types. Each has specific requirements detailed below.

\subsection{Option 1: Corpus Analysis}

\begin{graybox}
\faClipboardCheck\ \textbf{Requirements:}
\begin{itemize}
  \item Use an open source corpus or one from the Linguistic Data Consortium.
    \begin{itemize}
      \item UR students have free LDC access. See \href{https://www.library.rochester.edu/about/policies/linguistic-data-consortium-ldc}{this page}.
    \end{itemize}
  \item Develop data extraction tools.
    \begin{itemize}
      \item \textbf{Good:} I used spaCy to extract all sentences with passive constructions from the NOW corpus, then calculated their frequency by genre,~\ldots
      \item \textbf{Bad:} I searched COCA for what preceded this one word, then copied and pasted from the table it output into the paper.
    \end{itemize}
  \item Create processed dataset for analysis using best practices for data representation learned in the course---e.g. tidy data principles.
  \item Design analysis pipeline to test a specific linguistic hypothesis about patterns in the corpus.
\end{itemize}
\end{graybox}

\newpage

\subsection{Option 2: Judgment/Reading Experiment Design}

\begin{graybox}
\faFlask\ \textbf{Requirements:}
\begin{itemize}
  \item Design complete experimental protocol using experimental design principles from the course.
  \item Specify participant recruitment strategy and inclusion/exclusion criteria.
  \item Develop stimuli set for testing specific linguistic hypothesis.
  \item Define the procedure for data collection.
  \item Create a data collection instrument using a standard framework (PsychoPy, jsPsych, etc.).
  \item Generate simulated datasets.
    \begin{itemize}
      \item Unlikely you will be able to collect data from human subjects within the project timeline.
    \end{itemize}
  \item Develop a full analysis pipeline, including power analysis based on simulated data.
\end{itemize}
\end{graybox}

\subsection{Option 3: Mechanistic Interpretation Study}

\begin{graybox}
\faBrain\ \textbf{Requirements:}
\begin{itemize}
  \item Select set of pre-trained deep learning models of interest.
  \item Choose probing method(s) for mechanistic interpretation.
  \item Identify linguistic phenomenon(s) to investigate.
  \item Design probing experiments to:
    \begin{itemize}
      \item Isolate specific model components (layers, attention heads, neurons).
      \item Measure model behavior on controlled inputs.
    \end{itemize}
  \item Test hypothesis about internal workings of a deep learning model.
  \item Use statistical techniques developed in this course.
  \item Provide mechanistic interpretation of model behavior.
\end{itemize}
\end{graybox}

\section{Git Repository Structure}

All projects must be submitted via a Git repository with the following structure:

\newpage

\begin{bluebox}
\renewcommand{\arraystretch}{1.08}
\begin{tabularx}{\textwidth}{>{\ttfamily\scriptsize}p{0.49\textwidth}>{\scriptsize}X}
project-root/ & \\
\textbar-- README.md & Project overview and setup instructions\\
\textbar-- DESCRIPTION/pyproject.toml & Dependencies and environment setup\\
\textbar-- data/ & \\
\quad\textbar-- raw/ & Original data files\\
\quad\textbar-- processed/ & Cleaned/processed data\\
\textbar-- code/ & \\
\quad\textbar-- preprocessing/ & Data-preparation scripts\\
\quad\textbar-- analysis/ & Statistical-analysis scripts\\
\quad\textbar-- visualization/ & Plotting and figure generation\\
\textbar-- results/ & \\
\quad\textbar-- figures/ & Generated plots and visualizations\\
\quad\textbar-- tables/ & Statistical-output tables\\
\textbar-- paper/ & \\
\quad\textbar-- paper.pdf & Final paper submission\\
\quad\textbar-- paper.tex & LaTeX source (optional)\\
\textbar-- docs/ & Additional documentation using standard documentation tools (e.g. mkdocs)\\
\end{tabularx}
\end{bluebox}

\section{Paper Structure Requirements}

\subsection{All Project Types}

\begin{center}
\small
\renewcommand{\arraystretch}{1.18}
\begin{tabularx}{0.98\textwidth}{>{\bfseries}p{2.4cm}Xp{2.2cm}}
\toprule
\rowcolor{accentblue!10}
Section & Content Requirements & Page Range\\
\midrule
Abstract & Research question, methods, key findings & $\sim\frac{1}{3}$ page\\
Introduction & Background, motivation, hypothesis & 1--2 pages\\
Lit Review & Relevant prior work, theoretical framework & 2--3 pages\\
Methods & Detailed methodology, statistical approach & 2--3 pages\\
Results & Findings with figures/tables & 3--4 pages\\
Discussion & Interpretation, limitations, implications & 2--3 pages\\
Conclusion & Summary and future work & 0.5--1 page\\
References & Complete citations & 1--2 pages\\
\midrule
Total & & \textbf{12--18 pages}\\
\bottomrule
\end{tabularx}
\end{center}

\newpage

\subsection{Option-Specific Requirements}

\subsubsection{Option 1: Corpus Analysis}

\begin{graybox}
\begin{itemize}
  \item \textbf{Methods Section (2--3 pages):}
    \begin{itemize}
      \item Corpus selection justification with size, genre, and time period specifications.
      \item Data extraction pipeline architecture and tool specifications (e.g., spaCy, NLTK).
      \item Linguistic feature operationalization and annotation scheme.
      \item Statistical analysis plan including distributional assumptions and corrections for multiple comparisons.
    \end{itemize}
  \item \textbf{Results Section (3--4 pages):}
    \begin{itemize}
      \item Descriptive statistics of corpus composition and extracted features.
      \item Frequency distributions and co-occurrence patterns with confidence intervals.
      \item Statistical tests of linguistic hypotheses with effect sizes.
      \item Visualizations: frequency plots, heatmaps, network graphs of linguistic relationships.
    \end{itemize}
  \item \textbf{Required Appendices:}
    \begin{itemize}
      \item Sample of extracted data (minimum 100 examples) with annotations.
      \item Complete extraction scripts with documentation.
      \item Corpus preprocessing decisions and edge case handling.
      \item Inter-annotator agreement metrics if manual annotation was performed.
    \end{itemize}
\end{itemize}
\end{graybox}

\subsubsection{Option 2: Experiment Design}

\begin{graybox}
\begin{itemize}
  \item \textbf{Methods Section (2--3 pages):}
    \begin{itemize}
      \item Experimental design specification (within/between subjects, factorial structure, counterbalancing).
      \item Participant demographics, recruitment strategy, and inclusion/exclusion criteria.
      \item Stimuli construction methodology with linguistic controls and norming procedures.
      \item Power analysis with assumed effect sizes based on prior literature.
      \item Data collection procedure timeline and experimental interface details.
    \end{itemize}
  \item \textbf{Results Section (3--4 pages):}
    \begin{itemize}
      \item Simulated data generation process with parameter justification.
      \item Analysis of simulated datasets showing expected patterns.
      \item Power curves and sample size recommendations.
      \item Mixed-effects model specifications and random effects structure.
      \item Sensitivity analyses for different effect size scenarios.
    \end{itemize}
  \item \textbf{Required Appendices:}
    \begin{itemize}
      \item Complete stimuli list with item characteristics and counterbalancing scheme.
      \item Experimental interface screenshots and participant instructions.
      \item Simulation description and parameters for reproducibility.
    \end{itemize}
\end{itemize}
\end{graybox}

\newpage

\subsubsection{Option 3: Mechanistic Interpretation}

\begin{graybox}
\begin{itemize}
  \item \textbf{Methods Section (2--3 pages):}
    \begin{itemize}
      \item Model selection justification with architecture specifications.
      \item Probing methodology details (linear probes, attention analysis, causal interventions).
      \item Dataset construction for probing experiments with linguistic controls.
      \item Statistical framework for interpreting probe accuracy and significance.
      \item Comparison methodology across layers, attention heads, or model variants.
    \end{itemize}
  \item \textbf{Results Section (3--4 pages):}
    \begin{itemize}
      \item Layer-wise analysis of linguistic feature encoding.
      \item Attention pattern visualizations and interpretations.
      \item Statistical validation of mechanistic hypotheses.
      \item Comparison of representations across model architectures.
      \item Ablation study results showing component importance.
    \end{itemize}
  \item \textbf{Required Appendices:}
    \begin{itemize}
      \item Model specifications, hyperparameters, and training details.
      \item Complete probing datasets with linguistic annotations.
      \item Visualization gallery of attention patterns and activations.
      \item Reproducibility checklist with random seeds and compute requirements.
      \item Supplementary statistical analyses and robustness checks.
    \end{itemize}
\end{itemize}
\end{graybox}

\section{Grading Criteria}

\subsection{Presentation Assessment (20\% of Project Grade)}

\subsubsection{Required Slide Structure}

\begin{bluebox}
\textbf{Academic Presentation Format (15--20 slides total):}
\begin{enumerate}
  \item \textbf{Title Slide:} Title, authors, affiliations, date.
  \item \textbf{Outline/Roadmap:} Overview of talk structure.
  \item \textbf{Introduction/Motivation:} Why this research matters (2--3 slides).
  \item \textbf{Background/Literature Review:} Previous work and context (3--4 slides).
  \item \textbf{Research Questions:} Clear statement of hypotheses (1 slide).
  \item \textbf{Methods:} Data, tools, analysis approach (3--4 slides).
  \item \textbf{Results:} Key findings with visualizations (4--5 slides).
  \item \textbf{Discussion:} Interpretation and implications (2--3 slides).
  \item \textbf{Conclusions:} Summary and future work (1--2 slides).
  \item \textbf{Acknowledgments:} Collaborators, resources (1--2 slides).
  \item \textbf{References:} Key citations.
  \item \textbf{Appendix Slides:} Additional details for Q\&A.
\end{enumerate}
\end{bluebox}

\newpage

\subsubsection{Slide Design Requirements}

\begin{center}
\small
\renewcommand{\arraystretch}{1.18}
\begin{tabularx}{\textwidth}{>{\bfseries}p{2.6cm}Xp{1.3cm}}
\toprule
\rowcolor{accentblue!10}
Component & Criteria & Weight\\
\midrule
Visual Design & Clean layout, readable fonts (24pt+), appropriate color scheme & 15\%\\
Content Quality & Clear, concise text; effective use of figures/tables & 25\%\\
Structure & Logical flow, proper transitions, timing & 20\%\\
Delivery & Clear speech, eye contact, engagement with audience & 20\%\\
Q\&A Handling & Thoughtful responses, depth of understanding & 20\%\\
\bottomrule
\end{tabularx}
\end{center}

\subsubsection{Submission Requirements}

\begin{goldbox}
\faFilePdf\ \textbf{Presentation Materials to Submit:} PDF version of slides\\
\phantom{\faFilePdf\ }\textbf{Due:} 24 hours before presentation date
\end{goldbox}

\subsection{Code Assessment (20\% of Project Grade)}

\begin{graybox}
\faComments\ \textbf{Oral Assessment Format:}
\begin{itemize}
  \item 10--20 minute interview during finals period.
  \item Code will be shown without comments.
  \item Students must explain functionality line-by-line.
  \item Questions will test understanding of implementation choices.
\end{itemize}
\end{graybox}

\begin{center}
\small
\renewcommand{\arraystretch}{1.18}
\begin{tabularx}{\textwidth}{>{\bfseries}p{2.6cm}Xp{1.3cm}}
\toprule
\rowcolor{accentblue!10}
Component & Criteria & Weight\\
\midrule
Functionality & Code runs without errors, produces correct output & 40\%\\
Understanding & Can explain all code decisions during oral assessment & 30\%\\
Organization & Clear structure, appropriate modularization & 15\%\\
Documentation & README, inline documentation & 15\%\\
\bottomrule
\end{tabularx}
\end{center}

\newpage

\subsection{Paper Assessment (50\% of Project Grade)}

\begin{center}
\small
\renewcommand{\arraystretch}{1.18}
\begin{tabularx}{\textwidth}{>{\bfseries}p{2.6cm}Xp{1.3cm}}
\toprule
\rowcolor{accentblue!10}
Component & Criteria & Weight\\
\midrule
Research Question & Clear, testable hypothesis; appropriate scope & 15\%\\
Methodology & Appropriate statistical methods; sound design & 20\%\\
Analysis & Correct implementation; thorough exploration & 20\%\\
Results & Clear presentation; appropriate visualizations & 15\%\\
Discussion & Insightful interpretation; acknowledges limitations & 15\%\\
Writing Quality & Clear, concise, well-organized & 10\%\\
Citations & Complete, properly formatted & 5\%\\
\bottomrule
\end{tabularx}
\end{center}

\section{Detailed Grading Rubric}

\subsection{Exemplary (A: 93--100\%)}
\begin{bluebox}
\begin{itemize}
  \item \textbf{Code:} Elegant, efficient implementation with excellent documentation.
  \item \textbf{Analysis:} Sophisticated multifaceted statistical approach that uses advanced techniques appropriately.
  \item \textbf{Paper:} Publication-quality writing with novel insights.
  \item \textbf{Oral:} Demonstrates deep understanding, can discuss alternatives.
\end{itemize}
\end{bluebox}

\subsection{Proficient (B: 83--92\%)}
\begin{graybox}
\begin{itemize}
  \item \textbf{Code:} Correct implementation with good organization.
  \item \textbf{Analysis:} Appropriate methods correctly applied.
  \item \textbf{Paper:} Clear writing with sound conclusions.
  \item \textbf{Oral:} Can explain all code with minor prompting.
\end{itemize}
\end{graybox}

\subsection{Developing (C: 73--82\%)}
\begin{goldbox}
\begin{itemize}
  \item \textbf{Code:} Mostly functional with some errors or inefficiencies.
  \item \textbf{Analysis:} Basic methods with some misunderstandings.
  \item \textbf{Paper:} Adequate writing with unclear sections.
  \item \textbf{Oral:} Struggles to explain some code sections.
\end{itemize}
\end{goldbox}

\newpage

\subsection{Inadequate (D--F: \texorpdfstring{$<73\%$}{<73\%})}
\begin{redbox}
\begin{itemize}
  \item \textbf{Code:} Major errors or incomplete implementation.
  \item \textbf{Analysis:} Incorrect methods or misinterpretation.
  \item \textbf{Paper:} Poor organization or missing components.
  \item \textbf{Oral:} Cannot explain core functionality.
\end{itemize}
\end{redbox}

\section{Submission Requirements}

\begin{goldbox}
\faCalendar\ \textbf{Timeline:}
\begin{itemize}
  \item \textbf{Pre-proposal Meeting:} Scheduled on or before Oct. 16.
  \item \textbf{Project Proposal:} Due Oct. 28 (Week 9).
  \item \textbf{Project Presentations:} Dec. 7, 9, or 14 (Weeks 15--16).
  \item \textbf{Paper \& Code Submission:} Due Dec. 14.
  \item \textbf{Oral Assessments:} Scheduled during the Dec. 18--23 finals period.
\end{itemize}
\end{goldbox}

\subsection{Proposal Requirements (1--2 pages)}
\begin{enumerate}
  \item Research question and hypothesis.
  \item Chosen project option and justification.
  \item Data source or experimental-design overview.
  \item Proposed statistical methods.
  \item Timeline for completion.
\end{enumerate}

\subsection{Presentation Requirements (15 minutes)}
\begin{enumerate}
  \item Research question and motivation.
  \item Methodology overview.
  \item Preliminary results (if available).
  \item Challenges and next steps.
  \item 5 minutes for questions.
\end{enumerate}

\newpage

\section{Academic Integrity}

\begin{redbox}
\faExclamationTriangle\ \textbf{Important Policies:}
\begin{itemize}
  \item All code must include citations for external sources.
  \item LLM-generated code must be disclosed and you must provide prompts used.
  \item You may not use LLMs for writing the paper, except to generate tables, figures, etc.
\end{itemize}
\end{redbox}

\section{Resources and Support}

\subsection{Corpus Resources}

\subsubsection{Institutional \& Licensed Corpora}
\begin{itemize}
  \item \textbf{LDC:} Linguistic Data Consortium (\href{https://www.ldc.upenn.edu/}{ldc.upenn.edu})
  \item \textbf{BNC:} British National Corpus - 100M words balanced (\href{https://www.english-corpora.org/bnc/}{english-corpora.org/bnc})
\end{itemize}

\subsubsection{Large Web-Scale Corpora}
\begin{itemize}
  \item \textbf{Common Crawl:} 50+ billion web pages, petabytes of data (\href{https://commoncrawl.org/}{commoncrawl.org})
  \item \textbf{NOW Corpus:} 15.6+ billion words, updated monthly (\href{https://www.english-corpora.org/now/}{english-corpora.org/now})
  \item \textbf{GloWbE:} 1.9 billion words web-based English (\href{https://www.english-corpora.org/glowbe/}{english-corpora.org/glowbe})
  \item \textbf{COCA:} Corpus of Contemporary American English (\href{https://www.english-corpora.org/coca/}{english-corpora.org/coca})
\end{itemize}

\subsubsection{Multilingual \& Parallel Corpora}
\begin{itemize}
  \item \textbf{OPUS:} Largest parallel corpus, 100+ languages (\href{https://opus.nlpl.eu/}{opus.nlpl.eu})
  \item \textbf{Universal Dependencies:} Treebanks for 100+ languages (\href{https://universaldependencies.org/}{universaldependencies.org})
  \item \textbf{Tatoeba:} Community sentences in 400+ languages (\href{https://tatoeba.org/}{tatoeba.org})
  \item \textbf{OpenSubtitles:} Movie/TV subtitles, 60+ languages (\href{https://www.opensubtitles.org/}{opensubtitles.org})
\end{itemize}

\subsubsection{Specialized \& Literary Corpora}
\begin{itemize}
  \item \textbf{Project Gutenberg:} 60,000+ free books (\href{https://www.gutenberg.org/}{gutenberg.org})
  \item \textbf{Wikipedia Dumps:} Full Wikipedia text (\href{https://dumps.wikimedia.org/}{dumps.wikimedia.org})
  \item \textbf{Reddit Corpus:} Social media discussions (\href{https://pushshift.io/reddit/}{pushshift.io/reddit})
  \item \textbf{ArXiv:} Scientific papers, 270GB fulltext (\href{https://arxiv.org/help/bulk_data}{arxiv.org/bulk\_data})
\end{itemize}

\subsubsection{Curated Dataset Collections}
\begin{itemize}
  \item \textbf{GitHub NLP Datasets:} Comprehensive list (\href{https://github.com/niderhoff/nlp-datasets}{github.com/niderhoff/nlp-datasets})
\end{itemize}

\newpage

\begin{itemize}
  \item \textbf{Hugging Face Datasets:} 70,000+ datasets (\href{https://huggingface.co/datasets}{huggingface.co/datasets})
\end{itemize}

\subsection{NLP Tools}
\begin{itemize}
  \item \textbf{spaCy:} Industrial-strength NLP (\href{https://spacy.io/}{spacy.io})
  \item \textbf{Stanford CoreNLP:} Java-based toolkit (\href{https://stanfordnlp.github.io/CoreNLP/}{stanfordnlp.github.io/CoreNLP})
  \item \textbf{NLTK:} Natural Language Toolkit for Python (\href{https://www.nltk.org/}{nltk.org})
  \item \textbf{Stanza:} Neural NLP pipeline (\href{https://stanfordnlp.github.io/stanza/}{stanfordnlp.github.io/stanza})
\end{itemize}

\vfill
\begin{center}
\scriptsize\color{darkgray}
This rubric is subject to clarification. Updates will be announced on Zulip.\\
Last updated: August 25, 2026
\end{center}

\end{document}
